\section{Experimental Setup}
\label{sec:setup}

\subsection{Research Questions}
\textbf{RQ1 (Effectiveness)}: How does HSP-Agile compare with one-shot and iterative test-feedback baselines in strict success and strict formal conformance?\\
\textbf{RQ2 (Efficiency)}: What is the latency trade-off of HSP-Agile relative to iterative test-feedback under the same task set and run budget?\\
\textbf{RQ3 (Component Contribution)}: How much do formal checking, repair loop, and pattern guard each contribute to correctness and prevention metrics under hard-task conditions?

\subsection{Benchmark Construction}
The benchmark combines:
\begin{itemize}
  \item \textbf{Base tasks}: extracted from Agile-SOFL examples.
  \item \textbf{Hard synthetic tasks}: precedence-sensitive multi-output scenarios generated by a scripted task generator.
\end{itemize}
Each hard task contains five inputs, three outputs, and up to nine ordered scenarios.
To avoid trivial synthetic inflation, we enforce three constraints in generation: overlapping guard regions, non-uniform boundary values, and at least one scenario-order ambiguity risk.
All generated tasks are stored with deterministic seeds to enable exact regeneration.
Tasks that fail parser validity or scenario consistency checks are discarded before inclusion.

\subsection{Threat-Oriented Mutant Design}
We evaluate prevention under two mutant layers:
\begin{itemize}
  \item \textbf{Specification-confusion mutants}: inject faults at the requirement/specification interpretation level (e.g., condition polarity inversion, scenario-order swap, output binding mismatch).
  \item \textbf{Implementation-screening mutants}: inject faults in candidate implementations that may pass partial tests but violate strict scenario semantics or high-severity patterns.
\end{itemize}
Mutant operators are applied with fixed seeds and archived generation logs.
When an operator yields non-compilable artifacts, the mutant is excluded and counted in validity statistics rather than silently dropped.

\subsection{Compared Modes}
\begin{itemize}
  \item \texttt{B0}: reference implementation baseline.
  \item \texttt{B1}: one-shot LLM.
  \item \texttt{B2}: iterative test-feedback.
  \item \texttt{M}: full HSP-Agile (formal + repair + pattern).
  \item \texttt{A1/A2/A3}: ablations removing one key component.
\end{itemize}
All methods are executed under the same prompt template family and maximum-attempt budget unless otherwise specified.

\subsection{Execution Protocol}
Each method is run per task under a common attempt budget \(K\) and identical stopping policy.
For stochastic methods, we execute multiple seeded runs and aggregate at task level before cross-method significance testing.
Prompt templates and checker settings are fixed across compared modes to isolate the effect of methodological components instead of prompt engineering variance.
Any mode-specific preprocessing (e.g., pattern-screening enable/disable in ablations) is explicitly logged.

\subsection{Metrics}
Primary metrics:
\begin{itemize}
  \item strict success rate,
  \item strict formal conformance $\mathrm{Conf}(c,t)$ (Equation~\ref{eq:conformance}),
  \item mutation kill rate $\mu_k$,
  \item prevention detection rate (PDR),
  \item false accept rate (FAR),
  \item latency.
\end{itemize}
RQ1 is evaluated mainly by strict success and strict formal conformance; RQ2 by latency; and RQ3 by ablation deltas, PDR, and FAR.

\subsection{Statistics}
For paired per-task comparisons, we use the Wilcoxon signed-rank test for correctness-related outcomes.
For latency distributions, we use the Mann-Whitney U test.
We report effect sizes using rank-biserial correlation and provide adjusted \(p\)-values when multiple pairwise comparisons are conducted.
We avoid inferential claims when sample size assumptions are not met; such cases are presented descriptively.

\subsection{Reproducibility}
All results are produced through a scripted pipeline with deterministic seeds; raw execution logs are retained in a versioned artifact store.
Data preparation, aggregation, and plotting are separated into distinct pipeline stages to support independent reruns of each phase.
To keep reporting conservative, tables/figures include only completed reruns; pending reruns are marked \texttt{TBD} instead of estimated.
